% Generated draft: requires mathematical and source review.
\documentclass[11pt]{article}
\usepackage[margin=25mm]{geometry}
\usepackage{fontspec}
\setmainfont{Latin Modern Roman}
\usepackage{amsmath,amssymb,mathtools,unicode-math}
\setmathfont{Latin Modern Math}
\usepackage{enumitem,needspace,xurl,hyperref}
\setlength{\parindent}{0pt}
\setlength{\parskip}{5pt}
\setlength{\emergencystretch}{3em}
\newcommand{\N}{\mathbb{N}}
\newcommand{\Z}{\mathbb{Z}}
\newcommand{\Q}{\mathbb{Q}}
\newcommand{\R}{\mathbb{R}}
\newcommand{\C}{\mathbb{C}}
\newcommand{\F}{\mathbb{F}}
\newcommand{\A}{\mathbb{A}}
\newcommand{\PP}{\mathbb{P}}
\newcommand{\E}{\mathbb{E}}
\newcommand{\eps}{\varepsilon}
\newcommand{\dd}{\,\mathrm{d}}
\begin{document}
\section*{UnsolvedMath ID: 801014}
\textbf{Source problem code:} catalogue entry
\subsection{Definitions and mathematical statement}
The supplied record refers to a family of games \(G_k\), but does not define their rules, legal plays, information available to the players, or winning conditions. It also does not explicitly specify the domain of the index \(k\).

Let \(I\) denote the intended, unspecified index domain. For a fully specified game \(G_k\), let \(\Sigma_1(G_k)\) denote its set of legal strategies for Player 1, and let \(\operatorname{Win}_1(G_k,\sigma)\) mean that \(\sigma\) is a winning strategy for Player 1 under that game's rules. These are placeholders for definitions absent from the record.

The question can therefore be recorded only schematically as follows: is it true that
\[
\forall k\in I\;\bigl(k\geq 4\Longrightarrow \exists\sigma\in\Sigma_1(G_k)\;\operatorname{Win}_1(G_k,\sigma)\bigr)?
\]
The source reports that Player 1 has a winning strategy for \(k=2,3\), attributing this result to [Ba], and describes the cases \(k\geq4\) as unknown. Neither this reported result nor the supplied open status is independently verified here. A precise mathematical statement requires the missing definition of \(G_k\) and its index domain.
\subsection{Short English statement}
For each intended index \(k\geq4\), does Player 1 have a winning strategy in \(G_k\)? The game and index domain are not defined in the supplied record, so this remains a schematic statement requiring mathematical review.
\subsection{Sources}
\begin{itemize}
\item Local UnsolvedMath registry: permanent ID \(801014\), source problem code \texttt{catalogue entry}, titled \emph{Open Problems in Mathematical Logic --- Problem 26}. The supplied registry labels the problem open and attributes it to Barmpalias, Gao et al., with proposed year 2026.
\item Supplied source list: Barmpalias, Gao et al., \emph{Open Problems in Mathematical Logic} (2026), Problem 26, \url{https://arxiv.org/abs/2608.26628}. The record gives extraction method \texttt{arxiv\_latex} and access date 2026-10-16. These provenance details are supplied metadata; no source visit or independent verification is claimed.
\item [Ba]: the supplied statement attributes the winning-strategy result for \(k=2,3\) to this reference label. No corresponding author, title, publication details, or link is supplied, so this reference cannot be resolved from the record.
\end{itemize}
\end{document}
